% Options for packages loaded elsewhere
\PassOptionsToPackage{unicode}{hyperref}
\PassOptionsToPackage{hyphens}{url}
\documentclass[
  ignorenonframetext,
]{beamer}
\newif\ifbibliography
\usepackage{pgfpages}
\setbeamertemplate{caption}[numbered]
\setbeamertemplate{caption label separator}{: }
\setbeamercolor{caption name}{fg=normal text.fg}
\beamertemplatenavigationsymbolsempty
% remove section numbering
\setbeamertemplate{part page}{
  \centering
  \begin{beamercolorbox}[sep=16pt,center]{part title}
    \usebeamerfont{part title}\insertpart\par
  \end{beamercolorbox}
}
\setbeamertemplate{section page}{
  \centering
  \begin{beamercolorbox}[sep=12pt,center]{section title}
    \usebeamerfont{section title}\insertsection\par
  \end{beamercolorbox}
}
\setbeamertemplate{subsection page}{
  \centering
  \begin{beamercolorbox}[sep=8pt,center]{subsection title}
    \usebeamerfont{subsection title}\insertsubsection\par
  \end{beamercolorbox}
}
% Prevent slide breaks in the middle of a paragraph
\widowpenalties 1 10000
\raggedbottom
\AtBeginPart{
  \frame{\partpage}
}
\AtBeginSection{
  \ifbibliography
  \else
    \frame{\sectionpage}
  \fi
}
\AtBeginSubsection{
  \frame{\subsectionpage}
}
\usepackage{iftex}
\ifPDFTeX
  \usepackage[T1]{fontenc}
  \usepackage[utf8]{inputenc}
  \usepackage{textcomp} % provide euro and other symbols
\else % if luatex or xetex
  \usepackage{unicode-math} % this also loads fontspec
  \defaultfontfeatures{Scale=MatchLowercase}
  \defaultfontfeatures[\rmfamily]{Ligatures=TeX,Scale=1}
\fi
\usepackage{lmodern}
\ifPDFTeX\else
  % xetex/luatex font selection
\fi
% Use upquote if available, for straight quotes in verbatim environments
\IfFileExists{upquote.sty}{\usepackage{upquote}}{}
\IfFileExists{microtype.sty}{% use microtype if available
  \usepackage[]{microtype}
  \UseMicrotypeSet[protrusion]{basicmath} % disable protrusion for tt fonts
}{}
\makeatletter
\@ifundefined{KOMAClassName}{% if non-KOMA class
  \IfFileExists{parskip.sty}{%
    \usepackage{parskip}
  }{% else
    \setlength{\parindent}{0pt}
    \setlength{\parskip}{6pt plus 2pt minus 1pt}}
}{% if KOMA class
  \KOMAoptions{parskip=half}}
\makeatother
\setlength{\emergencystretch}{3em} % prevent overfull lines
\providecommand{\tightlist}{%
  \setlength{\itemsep}{0pt}\setlength{\parskip}{0pt}}
\usepackage{bookmark}
\IfFileExists{xurl.sty}{\usepackage{xurl}}{} % add URL line breaks if available
\urlstyle{same}
\hypersetup{
  hidelinks,
  pdfcreator={LaTeX via pandoc}}

\author{\texorpdfstring{}{}}
\date{}

\begin{document}

\begin{frame}{Abstract}
\protect\phantomsection\label{abstract}
This manuscript maps the structural collapse of British bimetallism
(1717--1844) onto the mythopoetic architecture of William Blake's
prophetic epics, arguing the disintegration of the gold-silver standard
materially enacted the cosmological fracture Blake dramatized. The
architecture of this entropy formalized on 21 September 1717, when Isaac
Newton, Master of the Royal Mint, established the guinea at 21 silver
shillings---a ratio of 15.21:1 that structurally overvalued gold against
continental markets. Under Gresham's Law, this metric disparity drove
English silver to France and the Far East, hollowed out the domestic
currency, and set the stage for the Bank Restriction crisis of February
1797. When the Bank of England's physical gold reserves collapsed
against a ten-fold circulation of paper notes, the British economy
underwent a systemic \emph{clinamen}---a swerve into the flux of fiat
abstraction. This groundlessness, extensively debated during the
Bullionist Controversy, culminated in the 1816 Coinage Act's statutory
establishment of the gold sovereign and the 1844 Bank Charter Act's
enclosure of value within the ledger. We trace this entropic trajectory
across the Atlantic to the American ``Crime of 1873'' and the 1896
``Cross of Gold'' crusade, revealing a transatlantic monetary framework
Blake first diagnosed in \emph{America A Prophecy}.

To counter this macro-economic atomization, Blake developed his
artisanal \emph{Fourfold Vision} as a systematic epistemology designed
to rescue human valuation from the mechanistic paradigms of capitalist
quantification. Blake's lifelong opposition to Newtonian ``Single
Vision'' aligns with his structural critique of the reductive logic of
monometallism. Where Newton's Mint fixed value within isolated metallic
quanta, Blake asserted the qualitative depths of imaginative perception.
His explicit marginalia against gold (``a guinea is more beautiful than
the sun'' to a miser), his condemnation of ``allegoric riches,'' and his
defense of Thomas Paine constitute a materialist indictment of the
commodification of human experience. By analyzing the physical Mint
alongside Los's spiritual forges as dual engines of socio-metaphysical
reproduction---the former quantifying labor into universal equivalents,
the latter transmuting material into the ``Human Form Divine'' through
relief-etching---this manuscript repositions Blake as a foundational
critic of financialized alienation. His prophetic stance bridges early
industrial monetary policy with the Marxist critique of commodity
fetishism, anticipating contemporary debates over fiat currencies
unmoored from their human ground.
\end{frame}

\end{document}
